% Role C / kangw24. Theory source: PR #81 q--g audit and production ECT
% pair construction. Keep q and g as implementation parameters; the scientific
% object is the realized pair-spacing intervention.
\subsection{Realized pair spacing is the intervention}
\label{sec:realized-spacing}

ECT constructs a coupled pair by sampling an online time $t$ and selecting an
earlier target time $r<t$.  Write the realized separation as
$\Delta=t-r$.  In the implementation studied here, curriculum parameter $q$,
stage $j$, and global multiplier $g$ jointly determine
\begin{align}
  r_q(t,j)
  &=\max\!\left\{0,
      t\left(1-\frac{n(t)}{q^{j+1}}\right)\right\},
  \\
  \Delta_{q,g}(t,j)
  &=\min\!\left\{t,\,g\bigl(t-r_q(t,j)\bigr)\right\},
  \qquad
  r_{q,g}(t,j)=t-\Delta_{q,g}(t,j),
  \label{eq:realized-spacing-construction}
\end{align}
where $n(t)$ is the fixed schedule-shape function.  The effective training
object is therefore the joint law
\begin{equation}
  \mathcal{P}_{q,g}
  =\operatorname{Law}\!\left(J,T,\Delta_{q,g}(T,J)\right)
  \label{eq:effective-spacing-law}
\end{equation}
under the frozen stage and time samplers.  Neither nominal $q$, nominal $g$,
nor the mean of $\Delta$ alone specifies this law.

At a fixed unclipped stage, multiplying the separation by $g$ can be written
as a reparameterization of $q$.  For reference parameters $(q_0,g_0)$, the
same pair coordinates are obtained with $g'=1$ and
$q_{\mathrm{eff}}(j)=q_0/g_0^{1/(j+1)}$.  This identity makes $g$ a controlled
implementation probe of realized spacing, rather than an independent
mechanism.  Across multiple stages or at a clipping boundary, equality
requires matching the full law $\mathcal{P}_{q,g}$ rather than one scalar
setting.

Realized spacing is a composite intervention in the evaluated ECT objective.
For online output $D_t=f_\theta(x_t,t)$ and detached target
$T_r=\operatorname{sg}[f_\theta(x_r,r)]$, the per-sample loss is
\begin{equation}
  \ell(\theta;t,r)
  =\frac{1}{\Delta}\,
    \rho\!\left(D_t-T_r\right).
  \label{eq:spacing-objective}
\end{equation}
Changing $\Delta$ moves the target endpoint $r$ and changes the explicit
inverse-spacing weight.  IA and IB instantiate this intervention by holding
the paired training configuration fixed and changing
\texttt{global\_gap\_scale} from $1.0$ to $1.1$.  Their comparison identifies
the effect of that realized-spacing intervention under the evaluated
construction.  An effect of $g$ independent of $q$ remains untested.
